\documentclass[10pt]{NSP1}
%%%%%%%%%%%%%%%%%%%%%%%%%%%%%%%%%%%%%%%%%%%%%%%%%%%%%%%%%%%%%%%%%%%%%%%%%%%%%%%%%%%%%%%%%%%%%%%%%%%%%%%%%%%%%%%%%%%%%%%%%%%%%%%%%%%%%%%%%%%%%%%%%%%%%%%%%%%%%%%%%%%%%%%%%%%%%%%%%%%%%%%%%%%%%%%%%%%%%%%%%%%%%%%%%%%%%%%%%%%%%%%%%%%%%%%%%%%%%%%%%%%%%%%%%%%%
\usepackage{url,floatflt}
\usepackage{helvet,times}
\usepackage{psfig,graphics}
\usepackage{mathptmx,amsmath,amssymb,bm}
\usepackage{float}
\usepackage[bf,hypcap]{caption}
\usepackage{hyperref}
\usepackage[dvipsnames,table,xcdraw]{xcolor}
\usepackage{tabularx}
\usepackage{booktabs}
\usepackage{bm}
\usepackage{afterpage}
\usepackage{algorithm}
\usepackage{algpseudocode} % For pseudocode formatting
\usepackage{longtable}


\setcounter{MaxMatrixCols}{10}
\addtolength{\oddsidemargin}{.5in}
\addtolength{\evensidemargin}{-.5in}

\tolerance=1
\emergencystretch=\maxdimen
\hyphenpenalty=10000
\hbadness=10000
\topmargin=0.00cm
\def\sm{\smallskip}
\def\no{\noindent}
\def\firstpage{59}
\setcounter{page}{\firstpage}
\def\thevol{14}
\def\thenumber{1}
\def\theyear{2026}


\begin{document}




\titlefigurecaption{{\large \bf \rm Journal of Statistics Applications \& Probability }\\ {\it\small An International Journal}}

\title{Implementation of Lean Six-Sigma to Improve Academic Achievement in Saad Al-Abdullah Academy in Kuwait}

\author{Rashed M. F. Al-Therwah\hyperlink{author1}{$^{1,2,*}$}, Mahmoud El-Sharief\hyperlink{author2}{$^2$}, Ibrahim Mohamed Hassab-Allah\hyperlink{author3}{$^2$} and Mahmoud Heshmat\hyperlink{author4}{$^2$}}
\institute{ Ministry of Interior, Kuwait \and
 Department of Mechanical Design and Production Engineering, Assiut University, Egypt}

\titlerunning{Implementation of Lean Six-Sigma...}
\authorrunning{Rashed M. et al.}

%corresponding author email
\mail{r.altherwah@gmail.com}

\received{15 Jul. 2024}
\revised{20 Sep. 2024}
\accepted{29 Nov. 2024}
\published{1 Jan. 2025}

\abstracttext{Through performance development and continuous improvement, educational institutions aim to accomplish their objectives efficiently and effectively. This study offers valuable insights for managing educational institutions, with a focus on contemporary management strategies. The study examines the benefits of applying Lean Six Sigma to enhance students' academic performance at Saad Al-Abdullah Academy in Kuwait. It focuses on factors that impact cadets' performance. To gather information on quality drivers about the following elements believed to affect academic performance, a survey was distributed to teachers and students: System, Vision and Mission, Curriculum Design and Alignment, Academic Environment, Student Personality-Related Factors, Teacher-Related Factors, and Teaching Methods and Approaches. All the identified criteria were covered by the questions included in the questionnaire. The findings were analyzed using the Six Sigma methodology (DMAIC) to identify the most important elements. The following important quality factors were determined to have the biggest effects on academic achievement based on the input that was collected and analyzed using the Pareto Chart for both teachers and students: Student Personality-Related Factors (active classroom participation improves performance, teacher quality and competence are critical to student success, and cultivating a growth mindset boosts student abilities through effort and practice), Academy System (the Academy's overall goals motivate high performance), Curriculum Design and Alignment (repetition in course content), and Teacher-Related Factors (expertise, pedagogical skills, and the ability to create a positive classroom environment significantly influence student achievement).}

\keywords{----------------} 



\maketitle

\section{Introduction}

Education is the cornerstone of progress and an essential pillar of national advancement when it is carried out effectively. The success of educational institutions, particularly those with large student populations, is hindered by several problems and obstacles, as recent research has revealed. These problems could jeopardize the institution's seamless operation and lead to unsatisfactory student outcomes if not promptly remedied \cite{R1}.


\end{document}
